\documentclass[12pt,a4paper]{article}

% ─── Packages ────────────────────────────────────────────────────────────────
\usepackage{amsmath, amssymb, amsthm}
\usepackage{geometry}
\usepackage{hyperref}
\usepackage{tikz}
\usepackage{pgfplots}
\pgfplotsset{compat=1.18}

\usetikzlibrary{
	arrows.meta, calc, positioning, shapes.geometric,
	decorations.pathreplacing, backgrounds, fit, matrix
}

\usepackage[dvipsnames]{xcolor}

\usepackage{tcolorbox}
\tcbuselibrary{theorems, skins, breakable}

\usepackage{booktabs}
\usepackage{array}
\usepackage{enumitem}
\usepackage{mathrsfs}
\usepackage{fancyhdr}
\usepackage{titlesec}
\usepackage{microtype}
\usepackage{setspace}

% ─── Page Geometry ────────────────────────────────────────────────────────────
\geometry{margin=2.5cm, top=3cm, bottom=3cm}
\setlength{\headheight}{14pt}
\setstretch{1.15}

% ─── Colours ──────────────────────────────────────────────────────────────────
\definecolor{axiomblue}{RGB}{30, 80, 160}
\definecolor{defgreen}{RGB}{20, 130, 80}
\definecolor{thmred}{RGB}{160, 30, 40}
\definecolor{notegray}{RGB}{90, 90, 90}
\definecolor{shadeblue}{RGB}{230, 240, 255}
\definecolor{shadegreen}{RGB}{225, 248, 235}
\definecolor{shadered}{RGB}{255, 235, 235}
\definecolor{shadeyellow}{RGB}{255, 252, 220}
\definecolor{purpleshade}{RGB}{240, 230, 255}
\definecolor{purpledark}{RGB}{100, 30, 160}
\definecolor{orangeshade}{RGB}{255, 243, 225}
\definecolor{orangedark}{RGB}{180, 90, 0}

% ─── Header / Footer ─────────────────────────────────────────────────────────
\pagestyle{fancy}
\fancyhf{}
\fancyhead[L]{\small\textit{Linear Algebra}}
\fancyhead[R]{\small\textit{Change of Basis}}
\fancyfoot[C]{\thepage}
\renewcommand{\headrulewidth}{0.4pt}

% ─── Section Formatting ───────────────────────────────────────────────────────
\titleformat{\section}{\large\bfseries\color{axiomblue}}{\thesection}{1em}{}[\titlerule]
\titleformat{\subsection}{\normalsize\bfseries\color{defgreen}}{\thesubsection}{1em}{}
\titleformat{\subsubsection}{\normalsize\bfseries\color{notegray}}{\thesubsubsection}{1em}{}

% ─── Theorem Environments ─────────────────────────────────────────────────────
\tcbset{
	defstyle/.style={
		enhanced, colback=shadegreen, colframe=defgreen,
		fonttitle=\bfseries, breakable
	},
	thmstyle/.style={
		enhanced, colback=shadered, colframe=thmred,
		fonttitle=\bfseries, breakable
	},
	axiomstyle/.style={
		enhanced, colback=shadeblue, colframe=axiomblue,
		fonttitle=\bfseries, breakable
	},
	notestyle/.style={
		enhanced, colback=shadeyellow, colframe=orange!70!black,
		fonttitle=\bfseries, breakable
	},
	propstyle/.style={
		enhanced, colback=purpleshade, colframe=purpledark,
		fonttitle=\bfseries, breakable
	},
	corstyle/.style={
		enhanced, colback=orangeshade, colframe=orangedark,
		fonttitle=\bfseries, breakable
	}
}

\newtcbtheorem[number within=section]{definition}{Definition}{defstyle}{def}
\newtcbtheorem[number within=section]{theorem}{Theorem}{thmstyle}{thm}
\newtcbtheorem[number within=section]{corollary}{Corollary}{corstyle}{cor}
\newtcbtheorem[number within=section]{proposition}{Proposition}{propstyle}{prop}
\newtcbtheorem[number within=section]{remark}{Remark}{notestyle}{rem}
\newtcbtheorem[number within=section]{example}{Example}{notestyle}{ex}

% ─── Math Commands ────────────────────────────────────────────────────────────
\newcommand{\R}{\mathbb{R}}
\newcommand{\vv}[1]{\mathbf{#1}}
\newcommand{\vzero}{\mathbf{0}}
\newcommand{\norm}[1]{\left\lVert #1 \right\rVert}
\DeclareMathOperator{\Span}{Span}
\DeclareMathOperator{\rank}{rank}
\DeclareMathOperator{\nullity}{nullity}

\renewcommand{\qedsymbol}{$\blacksquare$}

% ══════════════════════════════════════════════════════════════════════════════
\begin{document}
	
	% ══════════════════════════════════════════════════════════════════════════════
	\section{Introduction}
	
	Changing bases is a fundamental technique in linear algebra that allows us to simplify problems by representing vectors in more convenient coordinate systems. A change of basis corresponds to switching from one ordered basis of a vector space to another, and it is implemented using matrix multiplication.
	
	In $\R^2$, for example, instead of using the standard basis $\{\vv{e}_1, \vv{e}_2\}$, we may use a different basis that better suits the geometry or physics of a problem, such as a tangent–normal frame.
	
	% ══════════════════════════════════════════════════════════════════════════════
	\section{Coordinates in $\R^2$}
	
	\begin{definition}{Standard Coordinates}{stdcoords}
		Any vector $\vv{x} \in \R^2$ can be written as
		\[
		\vv{x} = x_1 \vv{e}_1 + x_2 \vv{e}_2,
		\]
		where $(x_1, x_2)^T$ are called the coordinates of $\vv{x}$ with respect to the standard basis $\{\vv{e}_1, \vv{e}_2\}$.
	\end{definition}
	
	\begin{definition}{Coordinates with respect to a basis}{coordsbasis}
		Let $\{\vv{y}, \vv{z}\}$ be a basis for $\R^2$. Then every vector $\vv{x} \in \R^2$ can be written uniquely as
		\[
		\vv{x} = \alpha \vv{y} + \beta \vv{z}.
		\]
		The vector $(\alpha, \beta)^T$ is called the coordinate vector of $\vv{x}$ with respect to the ordered basis $[\vv{y}, \vv{z}]$.
	\end{definition}
	
	\begin{remark}{Ordering of a basis}{ordering}
		The order of basis vectors matters. If the ordered basis is changed from $[\vv{y}, \vv{z}]$ to $[\vv{z}, \vv{y}]$, then the coordinate vector changes from $(\alpha, \beta)^T$ to $(\beta, \alpha)^T$.
	\end{remark}
	
	% ══════════════════════════════════════════════════════════════════════════════
	\section{Transition Matrices in $\R^2$}
	
	\subsection{Constructing a transition matrix}
	
	Suppose we are given a new basis
	\[
	\vv{u}_1 = \begin{pmatrix}3 \\ 2\end{pmatrix}, \quad
	\vv{u}_2 = \begin{pmatrix}1 \\ 1\end{pmatrix}.
	\]
	
	We form the matrix
	\[
	U = \begin{pmatrix}
		3 & 1 \\
		2 & 1
	\end{pmatrix}.
	\]
	
	\begin{definition}{Transition Matrix}{transmatrix}
		Let $U = (\vv{u}_1 \ \vv{u}_2)$ be the matrix whose columns are basis vectors expressed in the standard basis. Then:
		\[
		\vv{x} = U\vv{c}
		\]
		maps coordinates $\vv{c}$ in the $\{\vv{u}_1,\vv{u}_2\}$-basis to standard coordinates.
	\end{definition}
	
	\subsection{Changing from new basis to standard basis}
	
	If
	\[
	\vv{c} =
	\begin{pmatrix}
		c_1 \\
		c_2
	\end{pmatrix},
	\quad \text{then} \quad
	\vv{x} = c_1 \vv{u}_1 + c_2 \vv{u}_2.
	\]
	
	Expanding:
	\[
	\vv{x} =
	\begin{pmatrix}
		3c_1 + c_2 \\
		2c_1 + c_2
	\end{pmatrix}
	=
	\begin{pmatrix}
		3 & 1 \\
		2 & 1
	\end{pmatrix}
	\begin{pmatrix}
		c_1 \\
		c_2
	\end{pmatrix}.
	\]
	
	\subsection{Inverse transformation}
	
	To recover coordinates in the $\{\vv{u}_1,\vv{u}_2\}$ basis from standard coordinates:
	\[
	\vv{c} = U^{-1}\vv{x}.
	\]
	
	\begin{definition}{Inverse Transition Matrix}{invtrans}
		If $U$ is a transition matrix from a basis $B$ to the standard basis, then $U^{-1}$ is the transition matrix from the standard basis back to $B$.
	\end{definition}
	
	% ══════════════════════════════════════════════════════════════════════════════
	\section{General Change of Basis in Vector Spaces}
	
	\subsection{Coordinate vectors}
	
	\begin{definition}{Coordinate vector in a general vector space}{coordgen}
		Let $V$ be a vector space with ordered basis
		\[
		E = \{\vv{v}_1, \vv{v}_2, \dots, \vv{v}_n\}.
		\]
		Then every vector $\vv{v} \in V$ can be written uniquely as
		\[
		\vv{v} = c_1 \vv{v}_1 + c_2 \vv{v}_2 + \cdots + c_n \vv{v}_n.
		\]
		The vector
		\[
		[\vv{v}]_E = (c_1, c_2, \dots, c_n)^T
		\]
		is called the coordinate vector of $\vv{v}$ with respect to $E$.
	\end{definition}
	
	\begin{remark}{Interpretation}{interp}
		The coordinate vector $[\vv{v}]_E$ represents $\vv{v}$ in $\R^n$, even if $\vv{v}$ itself belongs to an abstract vector space $V$. This allows abstract vector problems to be converted into matrix problems.
	\end{remark}
	
	% ══════════════════════════════════════════════════════════════════════════════
	\section{Summary}
	
	\begin{tcolorbox}[colback=shadeyellow, colframe=orange!70!black, title=Summary, breakable]
		\begin{itemize}
			\item A change of basis converts vectors from one coordinate system to another.
			\item Coordinates depend on the ordering of the basis.
			\item Transition matrices map coordinates between bases.
			\item If $U$ is a transition matrix, then $U^{-1}$ reverses the change of basis.
			\item Every finite-dimensional vector space admits coordinate representations via ordered bases.
		\end{itemize}
	\end{tcolorbox}
	
\end{document}